\documentclass[10pt]{article}
\usepackage[margin=20mm]{geometry}
\usepackage{amsmath,amssymb,mathtools,fontspec,unicode-math,xurl,hyperref}
\IfFontExistsTF{DejaVu Serif}{\setmainfont{DejaVu Serif}}{\setmainfont{Latin Modern Roman}}
\setmathfont{latinmodern-math.otf}
\setlength{\parindent}{0pt}\setlength{\parskip}{6pt}\sloppy
\title{Technical review of universality for monotone cellular automata}
\author{AI-assisted technical review}
\date{}
\begin{document}
\maketitle
Abstract: The exact 72-page journal version is examined. An absolute support-coordinate bound in the proof of Lemma 6.12 is false for a translated family of admissible three-neighbour configurations. The counterexample does not contradict the lemma’s weighted-diameter conclusion. A replacement argument proving that conclusion is given, using the preceding geometric lemmas and an explicit choice of origin. The supported assessment is Minor. The principal probabilistic lower bound is not refuted; independent scientific adjudication remains pending.\par
\section*{Identity and scope}
The reviewed article is Paul Balister, Béla Bollobás, Robert Morris and Paul Smith, Universality for monotone cellular automata, J. Eur. Math. Soc. 28 (2026), 4725–4796, DOI 10.4171/JEMS/1666 (https://doi.org/10.4171/JEMS/1666) [paper]. The publisher PDF records receipt on 6 November 2022 and revision on 10 September 2024; the publisher landing page labels 3 June 2025 as Published. These dates and the 2026 issue year are distinct. The frozen Math2025 membership is unchanged. The arXiv metadata identifies 2203.13806v3, dated 7 November 2022; this report does not presume equality with that manuscript. Physical page \(k\) has printed page \(4724+k\).\par
All 72 pages were read as extracted text, including Appendix A and the references. Detailed checks followed the recursive resistance definition, the bounding-tree construction, closure on shifted lattices, iceberg containers, weighted diameter, subadditivity and the one-step probability induction. Decisive source formulas were also checked in actual page images. This coverage is not a claim to have independently reproved every lemma and every external reference.\par
For a fixed critical finite update family \(\mathcal U\) in \(d\) dimensions, Theorem 3.1, p. 4736, proves the lower bound \[p_c(\mathbb Z_n^d,\mathcal U)\ge
 \left(\frac1{\log_{(r-1)}n}\right)^{d-1+o(1)},
 \qquad r=r(\mathcal U),\quad n\longrightarrow\infty.\] Updates are finite nonempty subsets of \(\mathbb Z^d\setminus\{0\}\); the initial infections are Bernoulli with parameter \(p\). The constants may depend on the fixed update family and tree. The resistance lies in \(2,\ldots,d\) for critical families. Upper bounds and the subcritical part of the broader universality classification use companion papers; they are not all established within this article.\par
\section*{A confirmed local defect}
Problem 1 (Absolute support coordinates are not bounded by diameter).  In the proof of Lemma 6.12, p. 4767 (physical 43), equation (48) bounds each \(d_v(\mathcal X)\) by \(C\operatorname{wd}_{u_i}(G_{u_i}([K]_{u_j}))+O(1)\). The middle inequality of (49) similarly bounds \(d_v(\mathcal Y)\) by \(C\operatorname{wd}_{u_i}(\mathcal Y)\). Here Definition 6.2, p. 4761, defines \(d_v(B)=\sup_{x\in B}\langle x,v\rangle\), without centering. No origin is selected in this proof. Lemma 6.3 makes the weighted diameter translation invariant, but an individual support coordinate has no such invariance.\par
\subsection*{A complete admissible counterexample}
Take \(d=r=3\), \(\Lambda=\mathbb Z^3\), and the three-neighbour family \(\mathcal U=\mathcal N_3^3\): every three-element subset of \(\{\pm e_1,\pm e_2,\pm e_3\}\) is an update. Its stable set is the union of the three coordinate great circles. Choose the root children of the bounding tree to be \(\{\pm e_1,\pm e_2,\pm e_3\}\). For each such axis \(u\), choose the four children \[au+be_j,\quad au-be_j\qquad(e_j\perp u),\qquad
 a^2+b^2=1,\quad 0<b\ll a.\] This is an admissible two-level good bounding tree: the four local directions surround \(u\) on its small sphere, the axis has induced resistance two, each selected child is stable with the required induced resistance one, and every nonroot path is linearly independent. The radii can be made as small as required. For each path \(0\to u\to v\), the union of the two negative half-spaces contains at most two of the six nearest-neighbour vectors. It therefore contains no update, verifying the extra closure condition of Proposition 4.10, not merely individual stability. Choose the permitted root weights to be \(1/6\).\par
Let \(u_1=e_1\), \(k=1\), \(i=j=0\), and for integers \(N\ge1\) put \[K_N=\{(0,N,0),(0,N+1,1)\}.\] This finite nonempty set lies outside \(H_T(e_1)=\{x_1<0\}\); \(e_1\) is not a leaf. With only two seeds its ordinary closure is \([K_N]_0=K_N\). With the ocean \(x_1<0\) infected, both missing corners of the square receive one ocean neighbour and two planar neighbours. Consequently \[[K_N]_{e_1}=\{0\}\times\{N,N+1\}\times\{0,1\}\ne[K_N]_0.\] The displayed square is closed with the ocean: a site of the plane outside the square has at most one infected planar neighbour, and a site with \(x_1\ge1\) has at most one neighbour in the square. Thus the changed-closure case used in the proof really occurs.\par
Using precisely Definition 6.6, let \[A=G_0([K_N]_0)=\{0\}\times[N,N+1]\times[0,1].\] Then \(\operatorname{wd}_0(A)=1/3\). The set \(\mathcal Y=G_{e_1}(A)\) is given by \[\begin{gathered}
 x_1\ge0,\qquad ax_1+bx_2\le b(N+1),\qquad
 ax_1-bx_2\le-bN,\\
 ax_1+bx_3\le b,\qquad ax_1-bx_3\le0.
 \end{gathered}\] Its coordinate extrema give \[\mathcal X=G_0(\mathcal Y)
 =[0,b/(2a)]\times[N,N+1]\times[0,1].\] In particular, for \(v=e_2\in N_T^\to(0)\), \[d_{e_2}(\mathcal X)=d_{e_2}(\mathcal Y)=N+1.\] Equation (48) would give \(N+1\le C/3+O(1)\) for all \(N\), although its constants depend only on \(T\) and the update range. This is impossible. Likewise the middle inequality of (49) fails: \(\mathcal Y\) is simply translated by \(Ne_2\), so its weighted diameter stays fixed. All tree, closure, finiteness, lattice and nonleaf conditions are retained.\par
\subsection*{The required weighted estimate can be recovered}
The preceding counterexample concerns an unnecessarily strong intermediate statement. The actual conclusion (46), and hence Lemma 6.12, admits the following replacement proof.\par
Keep the notation of that proof and set \[S=[K]_{u_j},\qquad A=G_{u_i}(S),\qquad
 Y=G_{u_{j+1}}(A),\qquad X=G_{u_i}(Y).\] If the two closures agree, (46) is immediate. Otherwise \(A\) is strongly connected to the new half-space. If \(S\) already intersects it this is immediate; if not, the first new infection must use both the old closure and the new ocean, as in the paper’s closure argument. Also \(S\subset A\), since it stays outside the earlier oceans. Lemmas 6.10, 6.5 and 6.4 now yield \[\operatorname{diam}Y\le C\operatorname{wd}_{u_i}(A)+O(1),\qquad
 \operatorname{diam}A\le C\operatorname{wd}_{u_i}(A).                         \tag{1}\] There is no support-coordinate claim in (1).\par
If \(u_i=0\), \(G_0=F_0\) and the supporting constraints imply \(\operatorname{wd}_0(G_0(Y))=\operatorname{wd}_0(Y)\le\operatorname{diam}Y\), as required. Suppose \(u_i\ne0\), and choose \(z\in K\subset A\cap Y\). The set \(A=G_{u_i}(S)\) meets the plane \(\langle x,u_i\rangle=0\): starting from \(z\), move backwards parallel to \(u_i\) until the plane is reached. Every child normal \(v\in N_T^\to(u_i)\) has \(\langle v,u_i\rangle>0\), by Lemma 4.12 and the small-radius choice, so all its upper constraints remain satisfied. It follows that \[0\le\langle z,u_i\rangle\le\operatorname{diam}A.\] For each child normal, the definition of \(X=G_{u_i}(Y)\) gives \[d_v(X-z)\le d_v(Y-z)\le\operatorname{diam}Y.\] For the remaining normal, the truncation in \(G_{u_i}\) gives \[d_{-u_i}(X-z)\le\langle z,u_i\rangle\le\operatorname{diam}A.\] Use translation invariance of \(\operatorname{wd}_{u_i}\), positive weights and their sum one. These inequalities imply \[\operatorname{wd}_{u_i}(X)
 \le\max\{\operatorname{diam}A,\operatorname{diam}Y\}
 \le C\operatorname{wd}_{u_i}(A)+O(1).                          \tag{2}\] The containment (47) and monotonicity of weighted diameter turn (2) into (46). Iteration over the finite path proves Lemma 6.12. All constants still depend only on the fixed tree and update range.\par
\subsection*{Impact and repair boundary}
Lemma 6.12 feeds Lemma 6.14, which supplies the size estimates used in the Aizenman–Lebowitz and extremal lemmas and the one-step hierarchies. This is a central dependency, so a false line in its proof warrants explicit attention. However, (1)–(2) verify the exact weighted conclusion needed by those downstream uses. The counterexample refutes neither Lemma 6.12 nor the lower bound of Theorem 3.1. On this established scope the assessment is Minor: replace (48)–(49) with a centered argument. This grading rests on the demonstrated preservation of the needed conclusion, not simply on the apparent ease of editing the text.\par
\section*{Other checks, boundary cases and limitations}
The construction of a good tree was checked against the union-of-oceans closure condition; stability of each half-space alone would not suffice. Shifted lattices are explicitly allowed, and the geometric constants must not depend on their translation. The choice of iceberg type uses a permissible sequence and the deepest eligible vertex; replacing it by an arbitrary smallest droplet would change the proof. The root case \(u=0\), nonleaf restrictions, and the extra depth factor in \(\operatorname{diam}^{**}\) were retained in the probability induction.\par
Section 8 fixes \[(d-1)^{-1}>\varepsilon(2)>\cdots>\varepsilon(r)>0,\qquad
 x(s)=\exp_{(s-2)}(p^{-\varepsilon(s)}).\] It explicitly assumes \(p\) sufficiently small after the geometric constants have been fixed. The base estimate uses \(x^{d-1}p\to0\); the restrictions on \(\varepsilon(2)\) are essential. The disjoint-witness cases use the BK inequality; the crossing case uses events depending on disjoint subsets after the inner iceberg is fixed. Remark 7.7 makes the translated ocean a deterministic function of the selected droplets and vertex, which is needed in counting choices. These are not arbitrary independent events.\par
The final torus argument was checked for the scale condition \(n\gg x(r)\). For a lower bound it suffices, by monotonicity in \(p\), to work at \(p=(\log_{(r-1)}n)^{-1/\varepsilon}\), where \(0<\varepsilon<\varepsilon(r)\) is fixed. At that density, \[x(r)=\exp_{(r-2)}
 \bigl((\log_{(r-1)}n)^{\varepsilon(r)/\varepsilon}\bigr)\ll n,\] while \(x(r)\gg\log n\), so the final exponential bound defeats the polynomial number of droplets. The scale condition need not hold for every smaller \(p\); monotonicity transfers the conclusion from the displayed density. The iterated-exponential comparison and the polynomial count of faces in the final union bound were read.\par
The full chain of deterministic lemmas and all three probabilistic cases were read; the local geometric estimate above was independently reproved. Not every combinatorial subcase or external result has been independently reconstructed. The companion papers and all cited sharp-threshold results were not independently verified. A bounded search of publisher, arXiv and author sources did not locate a separate target erratum, supplement or implementation. Oxford ORA returned 403 and the Morris homepage returned empty redirect content. Absence from this search is not a defect, and no author code was run.\par
Dimension | State\par
Supported assessment | Minor: one false intermediate coordinate bound\par
Decisive evidence | An admissible translated bootstrap configuration; complete analytic geometry and a proved replacement estimate\par
Main-result status | No counterexample established; complete independent certification is not claimed\par
Scientific decision | Pending\par
Independent acceptance / human review | Not performed / none recorded\par
Author adoption | Unknown; not requested\par
1 P. Balister, B. Bollobás, R. Morris and P. Smith, Universality for monotone cellular automata, J. Eur. Math. Soc. 28 (2026), 4725–4796. https://doi.org/10.4171/JEMS/1666.\par
\section*{Supporting checks}
The independently written diagnostic `check\_0293.py` uses rational arithmetic to enumerate vertices of the displayed polytope and checks the finite local closure pattern. The analytic argument, not a finite truncation, proves the infinite-family contradiction. Source pages actually viewed were 1, 12, 26, 28, 37–39, 41–44, 57, 59, 65 and 66. All 72 pages were text-read.\par
Dispatcher UI observations and runtime self-description are separately attributed; no independent backend verification or scientific acceptance is asserted. The original acquisition token remains separate from the review token.\par
\end{document}
